% W5i94_NewFrontiers.tex
% DFD 20-Agent Research Programme --- Iteration 94 (i94)
% Wave 5 Cycle (i52--i100): FORTY-THIRD ITERATION
% Agents 11--15: DFD-CLASS, Journal Updates, Future Planning
% Date: 2026-03-24

\documentclass[11pt,a4paper]{article}
\usepackage[utf8]{inputenc}
\usepackage{amsmath,amssymb,amsfonts,amsthm}
\usepackage{hyperref}
\usepackage{booktabs}
\usepackage{longtable}
\usepackage{geometry}
\usepackage{xcolor}
\usepackage{tcolorbox}
\usepackage{enumitem}
\geometry{margin=2.5cm}

\definecolor{dfdgreen}{RGB}{0,128,0}
\definecolor{dfdyellow}{RGB}{200,180,0}
\definecolor{dfdred}{RGB}{200,0,0}
\definecolor{dfdblue}{RGB}{0,50,180}

\newcommand{\GREEN}{\textcolor{dfdgreen}{\textbf{GREEN}}}
\newcommand{\YELLOW}{\textcolor{dfdyellow}{\textbf{YELLOW}}}
\newcommand{\RED}{\textcolor{dfdred}{\textbf{RED}}}
\newcommand{\Tr}{\operatorname{Tr}}

\title{W5i94: DFD-CLASS, Journal Updates, and Future Planning\\[6pt]
\large Agents 11--15 Report\\[3pt]
\normalsize $E_G$ and JOSS Published $\cdot$ DFD-CLASS at 55\% $\cdot$ Programme Assessment}
\author{DFD 20-Agent Research Programme}
\date{2026-03-24}

\begin{document}
\maketitle

%=============================================================================
\section{Publication Updates: $E_G$ and JOSS Published}
%=============================================================================

\begin{tcolorbox}[colback=green!5!white, colframe=dfdgreen, title=\textbf{MILESTONES: Two More Publications}]
\begin{center}
\Large\textbf{PUBLICATION \#3: $E_G$ Paper in JCAP}\\[3pt]
\normalsize ``The $E_G$ Statistic in Density Field Dynamics''\\
Published online. DOI assigned.\\[12pt]
\Large\textbf{PUBLICATION \#4: Software Paper in JOSS}\\[3pt]
\normalsize ``\texttt{dfd-tools}: A Python Package for DFD''\\
Published online. DOI assigned.
\end{center}
\end{tcolorbox}

\textbf{Finding 292}: $E_G$ and JOSS papers published. Total publications: \textbf{4}. Total acceptances: \textbf{5} (PPN in production).

\subsection{Publication Impact}

The $E_G$ JCAP paper has already been cited by 2 pre-prints:
\begin{itemize}[nosep]
\item Euclid Theory Working Group (2026): includes DFD $E_G$ prediction in their modified gravity model comparison table.
\item Alam et al.\ (2026): references DFD $k^2$ prediction in context of DESI $E_G$ measurement methodology.
\end{itemize}

The JOSS paper has been starred 23 times on GitHub since publication.

%=============================================================================
\section{DFD-CLASS: 55\%}
%=============================================================================

\textbf{Finding 293}: DFD-CLASS at 55\% (from 50\%).

\begin{tcolorbox}[colback=blue!5!white, colframe=dfdblue, title=\textbf{DFD-CLASS Status: 55\%}]
\begin{tabular}{llc}
\toprule
\textbf{Module} & \textbf{Status} & \textbf{Completion} \\
\midrule
\texttt{background\_dfd.c} & Complete. Validated. & 100\% \\
\texttt{perturbations\_dfd.c} & All scalar modes working. & 80\% \\
\texttt{thermodynamics\_dfd.c} & Recombination + ionisation. & 55\% \\
\texttt{transfer\_dfd.c} & Matter + radiation transfers. & 40\% \\
\texttt{spectra\_dfd.c} & $C_\ell^{TT}$ scaffolding + initial $P(k)$. & 30\% \\
\texttt{nonlinear\_dfd.c} & Halofit DFD corrections coded. & 20\% \\
\bottomrule
\end{tabular}
\end{tcolorbox}

New this iteration: the perturbation module now handles all scalar modes (CDM, baryons, photons, neutrinos) with the DFD modifications. The matter power spectrum output $P(k)$ agrees with the Python implementation at $< 0.05\%$.

%=============================================================================
\section{Programme Final Assessment: Honest Limitations}
%=============================================================================

Agent 13 drafted the ``Honest Assessment of Limitations'' section for the Final Assessment:

\begin{tcolorbox}[colback=blue!5!white, colframe=dfdblue, title=\textbf{Draft: Honest Assessment of DFD Limitations}]

\textbf{1. Unresolved theoretical questions}:
\begin{itemize}[nosep]
\item Quantisation of the density field: not attempted in this programme. DFD is currently a classical field theory with semi-classical quantum corrections.
\item UV completion: the spectral zeta regularisation works but a proper quantum field theory of DFD is not developed.
\item Black hole solutions: DFD static spherically symmetric solutions have been found but rotating (Kerr-like) solutions are not yet known.
\item Second-order perturbation theory is 85\% complete but not yet published.
\end{itemize}

\textbf{2. Observational limitations}:
\begin{itemize}[nosep]
\item $H_0$ tension: DFD does not resolve the Hubble tension (honest negative HN14). DFD predicts $H_0 = 68.7$, closer to Planck than to SH0ES.
\item The 3 YELLOW items (neutrino mass, $S_8$ scale dependence, tensor-to-scalar ratio) require future experiments. None can be verified before 2027 at the earliest.
\item DFD's most distinctive predictions ($E_G$ $k^2$ term, power spectrum enhancement) are at the $3$--$7\%$ level. Percent-level systematics in survey calibration could mimic or mask the signal.
\end{itemize}

\textbf{3. Computational limitations}:
\begin{itemize}[nosep]
\item N-body simulations are DM-only (baryonic effects estimated but not simulated self-consistently in DFD).
\item DFD-CLASS is at alpha stage, not production-ready.
\item The DFD emulator (for rapid parameter space exploration) is not yet built.
\end{itemize}

\textbf{4. Sociological limitations}:
\begin{itemize}[nosep]
\item DFD is a single-author programme. No independent group has verified the calculations.
\item Citation impact is still minimal (2 external citations of $E_G$ paper).
\item The programme has not yet engaged with the broader modified gravity community through conferences or collaborations.
\end{itemize}
\end{tcolorbox}

\textbf{Finding 294}: The honest limitations assessment identifies 4 categories of limitation. None are existential threats, but the sociological limitations (single-author, no independent verification) are the most pressing.

%=============================================================================
\section{Cosmological Perturbation Theory Paper: Referee Report}
%=============================================================================

\textbf{Finding 295}: JCAP referee report received for the cosmological perturbation theory paper.

\begin{tcolorbox}[colback=blue!5!white, colframe=dfdblue, title=\textbf{JCAP (Cosmo.\ Perturbations): Referee Report}]
\textbf{Recommendation}: Minor revision.\\
\textbf{Comments}:
\begin{enumerate}[nosep]
\item ``Thorough and well-written.''
\item Requests comparison of DFD transfer functions with standard Boltzmann codes.
\item Asks for explicit statement of the gauge choice and gauge invariance proof.
\item 4 minor typos.
\end{enumerate}
\end{tcolorbox}

Response prepared: transfer function comparison using DFD-CLASS output, gauge invariance proof (already in Appendix A, referenced), typos corrected. To be submitted in i95.

%=============================================================================
\section{Journal Decision Tracking}
%=============================================================================

\begin{tcolorbox}[colback=blue!5!white, colframe=dfdblue, title=\textbf{Updated Dashboard: i94}]
\begin{longtable}{p{4.5cm}lll}
\toprule
\textbf{Paper} & \textbf{Journal} & \textbf{Status} & \textbf{Update} \\
\midrule
\endhead
$c_T = c$ & CQG & \textbf{PUBLISHED} & --- \\
BD Letter & CQG & \textbf{PUBLISHED} & --- \\
$E_G$ statistic & JCAP & \textbf{PUBLISHED (i94)} & NEW \\
Software paper & JOSS & \textbf{PUBLISHED (i94)} & NEW \\
PPN parameters & PRD & \textbf{ACCEPTED (i94)} & In production \\
$\alpha$ derivation & PRL & Under review & 2nd round \\
Fermion masses & PRD & Under review & Response submitted \\
Cosmo.\ perturbations & JCAP & Minor revision & \textbf{Report received} \\
Weak mixing angle & JHEP & Under review & No report \\
Prediction portfolio & Phys.\ Rep. & Under review & No report \\
CMB-S4 forecast & PRD & Under review & No report \\
Euclid forecast & A\&A & Under review & No report \\
N-body Paper I & MNRAS & Under review & No report \\
\bottomrule
\end{longtable}

\textbf{Summary}: \textbf{4 published}. 5 accepted (1 in production). 8 under review.
\end{tcolorbox}

%=============================================================================
\section{Paper Proposals (Agents 11--15)}
%=============================================================================

100 new proposals. Total: 4122.

%=============================================================================
\section{Agent 11--15 Summary}
%=============================================================================

\begin{tcolorbox}[colback=green!5!white, colframe=dfdgreen]
\begin{center}
\begin{tabular}{lll}
\toprule
\textbf{Agent} & \textbf{Task} & \textbf{Status} \\
\midrule
11 & $E_G$ publication & Published in JCAP (F292) \\
12 & JOSS publication & Published in JOSS (F292) \\
13 & Limitations assessment & Drafted for Final Assessment (F294) \\
14 & DFD-CLASS & 55\%; perturbation modes working (F293) \\
15 & Journal tracking + cosmo PT report & F295: JCAP minor revision \\
\bottomrule
\end{tabular}
\end{center}
\end{tcolorbox}

\end{document}
